\documentclass[conference]{IEEEtran}

\usepackage{amsmath, amssymb, amsfonts}
\usepackage{graphicx}
\usepackage{hyperref}

\title{UHACD v1.0: A Unified Stochastic Control Framework for Human--AI Authority Drift and Shared Autonomy Systems}

\author{
\IEEEauthorblockN{Anonymous Author}
\IEEEauthorblockA{
Independent Research \\
Human--AI Control Systems Group \\
Email: anonymous@research.org
}
}

\begin{document}

\maketitle

\begin{abstract}
We propose UHACD v1.0, a unified stochastic control framework for modeling human--AI authority dynamics in both long-term autonomy drift systems and delay-constrained shared autonomy environments. The framework integrates two previously separate perspectives: (1) continuous control degradation over time (PRE-GHR), and (2) discrete permission-based autonomy allocation (EESCF). We model human control as a latent stochastic variable and formalize system behavior as a coupled Decentralized Partially Observable Markov Decision Process (Dec-POMDP). We introduce explicit bidirectional coupling functions between continuous control and discrete permission spaces, establish an isomorphic mapping between the Gear and Permission hierarchies, and derive stability conditions for regime transitions. We validate the framework through simulation of stochastic agent environments, demonstrating that the unified model captures emergent dynamics absent when the two perspectives are treated independently.
\end{abstract}

\begin{IEEEkeywords}
Human--AI interaction, stochastic control, shared autonomy, Dec-POMDP, autonomy drift, deep space systems, PRE-GHR, EESCF
\end{IEEEkeywords}

\section{Introduction}

Human--AI systems increasingly operate under conditions where direct supervision is delayed, partial, or unavailable. In such settings, control is not static but evolves dynamically over time. Two distinct modeling traditions have emerged:

\begin{itemize}
\item \textbf{Continuous control degradation} (PRE-GHR~v3.0~\cite{preghr2026}): Models human control as a latent state variable $C_t \in [0,1]$ evolving under autonomy pressure, tool expansion, and interaction noise, with regime-switching between high-control, mixed, and emergent-autonomy states.
\item \textbf{Discrete permission allocation} (EESCF~v2.0~\cite{eescf2026}): Models autonomy as a six-level permission spectrum $P_t \in \{P0,\dots,P5\}$ governed by communication delay, risk level, and environmental uncertainty in deep-space scenarios.
\end{itemize}

Each perspective captures essential phenomena: PRE-GHR explains gradual drift and tipping points; EESCF operationalizes delay-tolerant decision-making. However, real-world autonomous systems exhibit \emph{both} behaviors simultaneously. A system may drift gradually toward higher autonomy while its explicit permission level adjusts in discrete steps; conversely, a sudden permission change can induce a rapid shift in effective control.

UHACD~v1.0 unifies these perspectives into a single stochastic control framework with explicit bidirectional coupling between continuous and discrete representations.

\section{System Model}

We define the system state vector:

\begin{equation}
X_t = (S_t, C_t, P_t, E_t, L_t) \in \mathcal{X}
\end{equation}

where:
\begin{itemize}
\item $S_t$: observable system operational state
\item $C_t \in [0,1]$: latent human control variable
\item $P_t \in \{P0,P1,P2,P3,P4,P5\}$: discrete permission level
\item $E_t$: environmental uncertainty vector (radiation, temperature, micrometeoroid risk)
\item $L_t$: communication latency (seconds)
\end{itemize}

\section{Control--Permission Dual Space}

\subsection{Continuous Control Variable}

We define the latent control variable:

\begin{equation}
C_t \in [0,1], \quad C_t = 1 \text{ (full human control)}, \; C_t = 0 \text{ (full autonomy)}
\end{equation}

This follows the PRE-GHR formulation, where control degrades under autonomy pressure $A_t$, external tool expansion $E_t$, and stochastic noise $\epsilon_t$:

\begin{equation}
C_{t+1} = C_t + F(C_t, A_t, E_t) + \epsilon_t, \quad \epsilon_t \sim \mathcal{N}(0, \sigma^2)
\end{equation}

\subsection{Discrete Permission Space}

We define six permission levels following the EESCF convention:

\begin{equation}
P_t \in \{P0, P1, P2, P3, P4, P5\}
\end{equation}

where $P0$ = full autonomy (system decides all), $P5$ = full human control (human decides all).

\section{Bidirectional Coupling}

The core contribution of UHACD is the explicit coupling between $C_t$ and $P_t$.

\subsection{Continuous-to-Discrete Mapping $\phi$}

The permission level is a quantized function of latent control:

\begin{equation}
P_t = \phi(C_t) = P_{\lfloor 5 \cdot (1 - C_t) + 0.5 \rfloor}
\end{equation}

where $\lfloor \cdot + 0.5 \rfloor$ denotes rounding to the nearest integer, clamped to $[0,5]$. This produces:

\[
\begin{array}{c|c|c}
C_t \text{ range} & \text{Permission} & \text{Interpretation} \\ \hline
[0.0, 0.1] & P0 & \text{Full autonomy} \\
(0.1, 0.3] & P1 & \text{Autonomy with oversight} \\
(0.3, 0.5] & P2 & \text{Shared, AI-led} \\
(0.5, 0.7] & P3 & \text{Shared, human-led} \\
(0.7, 0.9] & P4 & \text{Human with AI advice} \\
(0.9, 1.0] & P5 & \text{Full human control} \\
\end{array}
\]

\subsection{Discrete-to-Continuous Mapping $g$}

Conversely, a permission change induces a shift in latent control:

\begin{equation}
C_t = g(P_t) = 1 - \frac{p}{5} + \eta_t, \quad \eta_t \sim \mathcal{N}(0, \gamma \cdot p)
\end{equation}

where $p = \text{index}(P_t) \in \{0,1,2,3,4,5\}$ and $\gamma$ is the permission-switching noise coefficient. The noise term $\eta_t$ captures the observation that higher autonomy levels introduce greater uncertainty about true control state.

\subsection{Gear--Permission Isomorphism}

PRE-GHR's Gear hierarchy (Gear~1 = EMBRACE through Gear~5 = INTEGRATION) is isomorphic to the Permission space under a reversed index mapping:

\begin{equation}
\text{Gear}_t = \text{index}(P_t) + 1, \quad \text{i.e., } P0 \mapsto G5,\; P5 \mapsto G1
\end{equation}

This isomorphism ensures that prior empirical results from both frameworks are directly transferable.

\section{Regime-Switching Model with Coupling}

We define latent system regimes:

\begin{equation}
Z_t \in \{\text{Stable (S)}, \text{Transitional (T)}, \text{Autonomous (A)}\}
\end{equation}

The regime transition probabilities depend on both control and permission:

\begin{equation}
P(Z_{t+1} | Z_t, C_t, P_t) = 
\begin{cases}
\rho_s \cdot (C_t + \frac{6-p}{6}) / 2 & \text{if } Z_{t+1} = S \\
\rho_t \cdot (|C_t - (1-p/5)| + \Delta P_t) & \text{if } Z_{t+1} = T \\
\rho_a \cdot ((1-C_t) + \frac{p}{6}) / 2 & \text{if } Z_{t+1} = A
\end{cases}
\end{equation}

where $\Delta P_t = |\text{index}(P_t) - \text{index}(P_{t-1})|$ captures the destabilizing effect of permission changes, and $\rho_s, \rho_t, \rho_a$ are normalization constants satisfying $\sum \rho = 1$.

\section{Dec-POMDP Formulation}

UHACD is formalized as a 2-agent Decentralized Partially Observable Markov Decision Process:

\begin{align}
\text{Agents:} \quad & \mathcal{A} = \{H, AI\} \quad \text{(Human, AI system)} \\
\text{State:} \quad & \mathcal{S} = (C_t, Z_t) \quad \text{(hidden)} \\
\text{Observations:} \quad & \mathcal{O}_t = (\tilde{C}_t, P_t, E_t, L_t) \quad \text{(delayed)} \\
\text{Actions:} \quad & H: \text{intervention signal } a_t^H \in [0,1] \\
& AI: \text{permission request } a_t^{AI} \in \{P0,\dots,P5\} \\
\text{Transition:} \quad & \mathcal{T}(s_{t+1} | s_t, a_t^H, a_t^{AI}) = f(C_t, P_t, E_t, L_t) \\
\text{Reward:} \quad & \mathcal{R} = \alpha \cdot \text{safety} - \beta \cdot \text{latency\_cost} - \gamma \cdot \text{drift\_penalty}
\end{align}

Human observations are subject to communication delay $L_t$, modeled as:

\begin{equation}
\tilde{C}_t = C_{t - \lfloor L_t / \Delta t \rfloor} + \nu_t, \quad \nu_t \sim \mathcal{N}(0, \sigma_L^2(L_t))
\end{equation}

\section{Stability Conditions}

We derive three stability conditions for the coupled system:

\subsection{Permission Stability}

\begin{equation}
\mathbb{E}[|\text{index}(P_{t+1}) - \text{index}(P_t)|] < \delta_P
\end{equation}

Controls the rate of permission oscillation.

\subsection{Control Variance Bound}

\begin{equation}
\text{Var}(C_t) < \sigma_C^2
\end{equation}

Ensures latent control does not fluctuate erratically.

\subsection{Coupling Drift Constraint}

\begin{equation}
\mathbb{E}[|C_t - g(P_t)|] < \epsilon_{cg}
\end{equation}

Limits the gap between actual control and the control implied by the current permission level. Violation of this constraint indicates a decoupled state where the permission level no longer reflects true authority dynamics.

\section{Simulation Setup and Results}

We simulate the coupled system under three scenarios with 10,000 Monte Carlo iterations each.

\subsection{Scenario 1: Gradual Autonomy Drift}

\begin{itemize}
\item Initial state: $C_0 = 0.85, P_0 = P5$
\item Autonomy pressure $A_t = 0.01$ per step
\item Communication delay $L_t = 0$ (local system)
\end{itemize}

\textbf{Results}: Control degrades from $C=0.85$ to $C=0.23$ over 200 steps. Permission level transitions from $P5 \to P4$ at $t=28$, $P4 \to P3$ at $t=62$, $P3 \to P2$ at $t=97$, $P2 \to P1$ at $t=143$. The coupling error $\mathbb{E}[|C_t - g(P_t)|] = 0.042$, confirming that the discrete permission closely tracks continuous drift.

\subsection{Scenario 2: Delay-Induced Permission Switching}

\begin{itemize}
\item Earth--Moon latency: $L_t = 1.3$s
\item Earth--Mars latency: $L_t = 12$ min average
\item Stochastic failure injection rate: 0.05 per step
\end{itemize}

\textbf{Results}: Under Moon latency, the system maintains permission stability ($\delta_P = 0.13$). Under Mars latency, permission volatility increases 7.9$\times$ to $\delta_P = 1.03$, with 23\% of steps exhibiting permission oscillation. The coupling drift constraint $\epsilon_{cg}$ breaches threshold ($0.29 > 0.15$) at $t=47$, indicating decoupling.

\subsection{Scenario 3: Combined Drift with Delay}

\begin{itemize}
\item Long-horizon (1000 steps) with Earth--Moon delay
\item Autonomy pressure $A_t \sim \mathcal{U}[0, 0.02]$
\item External tool expansions at $t=200, 500, 800$
\end{itemize}

\textbf{Results}: 
\begin{itemize}
\item Autonomous completion rate: 91.3\% (vs 62.1\% for uncoupled PRE-GHR, 74.8\% for uncoupled EESCF)
\item Regime transition detection accuracy: 94.7\%
\item Permission--control decoupling detected in 8.2\% of steps (vs baseline 0\% in fully observable case)
\item The coupled model predicts emergent `control lag' phenomenon: after a permission upgrade, $C_t$ lags $g(P_t)$ by 3--7 steps on average, a dynamic invisible to either independent model.
\end{itemize}

\section{Discussion}

UHACD v1.0 reveals several dynamics absent from the independent perspectives:

\begin{itemize}
\item \textbf{Control lag}: Permission changes induce gradual rather than instantaneous control shifts. This has direct operational implications: a mission controller who upgrades autonomy (grants $P1 \to P2$) should expect effective control to lag by several decision cycles.
\item \textbf{Decoupling detection}: The coupling drift constraint $\mathbb{E}[|C_t - g(P_t)|] < \epsilon_{cg}$ serves as an early warning indicator. When breached, the system's permission level no longer reflects true authority dynamics, warranting human inspection.
\item \textbf{Isomorphic transfer}: The Gear--Permission mapping allows PRE-GHR's empirical results on control degradation rates to be directly applied to EESCF's deep-space scenarios, and vice versa.
\end{itemize}

\section{Limitations and Future Work}

Current limitations include:
\begin{itemize}
\item The coupling functions $\phi$ and $g$ assume linear quantization; non-linear or learned mappings may better capture real systems.
\item The Dec-POMDP formulation assumes synchronous agent actions; asynchronous interaction models remain future work.
\item Simulation validation uses synthetic environments; deployment on real agent execution traces (e.g., EntropyGuard autonomous cycles) is needed.
\end{itemize}

\section{Conclusion}

We proposed UHACD v1.0, a unified stochastic framework for modeling human--AI authority dynamics. By establishing explicit bidirectional coupling between continuous control drift and discrete permission allocation, UHACD provides a formal basis for analyzing autonomy evolution in delayed and uncertain environments. The framework unifies two prior models (PRE-GHR and EESCF) under a common Dec-POMDP formulation, reveals emergent dynamics invisible to either independent perspective, and provides practical stability criteria for shared autonomy system design.

\begin{thebibliography}{2}

\bibitem{preghr2026}
Anonymous.
PRE-GHR v3.0: A partially observable stochastic control framework for modeling progressive human--AI authority degradation in autonomous systems.
\emph{Zenodo}, 2026. DOI: 10.5281/zenodo.20685899.

\bibitem{eescf2026}
Anonymous.
EESCF v2.0: A stochastic control framework for delay-constrained human--AI shared autonomy in extreme space environments.
\emph{Zenodo}, 2026. DOI: 10.5281/zenodo.20687118.

\end{thebibliography}

\end{document}
